\clearpage\thispagestyle{empty}
\pdfbookmark[0]{Quellen zu Pallas (Dezember 1810)}{pallas1810dec-sources}
\begin{center}\textbf{Quellen zu Pallas (Dezember 1810)}\end{center}
\noindent Diese KI-Übertragung ins heutige Deutsch folgt dem Bericht in \href{https://gdz.sub.uni-goettingen.de/id/PPN236007467}{Gauß’ \textit{Werken}, Bd.~VI} (Göttingen, 1874), S.~320--322. Die entsprechenden Abschnitte erschienen in Beitrag LXII der Zeitschrift \href{https://archive.org/download/bub_gb_YCI4AAAAMAAJ/page/n4.jpg}{\textit{Monatliche Correspondenz}, Bd.~XXII} vom Dezember 1810, S.~\href{https://archive.org/download/bub_gb_YCI4AAAAMAAJ/page/n594.jpg}{591}, \href{https://archive.org/download/bub_gb_YCI4AAAAMAAJ/page/n595.jpg}{592}, \href{https://archive.org/download/bub_gb_YCI4AAAAMAAJ/page/n596.jpg}{593}, \href{https://archive.org/download/bub_gb_YCI4AAAAMAAJ/page/n597.jpg}{594} und \href{https://archive.org/download/bub_gb_YCI4AAAAMAAJ/page/n598.jpg}{595}. Das Exemplar der University of Michigan wurde von Google digitalisiert und ist im Internet Archive zugänglich.
\par\smallskip\noindent Die Zeitschrift leitet einen Brief von Gauß vom 2. Dezember ein. Die Werkausgabe setzt mit dem Zitat ein, enthält seine Vorhersage für den 15. Dezember und schließt mit seiner Antwort auf den Beobachtungsbericht. In der Tabelle auf S.~594 stehen in der Zeitschrift das Jahr 1810 und die Deklination $21^\circ\;58''7$. Die Werkausgabe gibt 1809 und $22^\circ\;58''7$ an. Bei den Bahnelementen auf S.~592 unterscheiden sich auch die Winkelminuten in der Zeile zur Exzentrizität und die Winkelsekunden der Neigung. Die Übertragung folgt der Werkausgabe. KI-Übertragung und KI-Satz.
\clearpage
